\title{Controlling Text-to-Image Diffusion by Orthogonal Finetuning}

\begin{abstract}
Text-to-image diffusion models have demonstrated remarkable proficiency in translating textual descriptions into photorealistic visual content. However, the challenge of steering these powerful generative models to accomplish specialized downstream objectives remains an open research problem. To address this, we present a systematic finetuning strategy called Orthogonal Finetuning (OFT), specifically designed for adapting text-to-image diffusion models for new tasks. Distinct from prior approaches, OFT is theoretically guaranteed to maintain hyperspherical energy, an important quantity reflecting the pairwise associations among neurons on the unit hypersphere. Our investigations reveal that retaining this energy is pivotal to sustaining the semantic generative capacity of text-to-image diffusion architectures. To further enhance the robustness of finetuning, we propose Constrained Orthogonal Finetuning (COFT), which adds a radius constraint within the hyperspherical manifold. Our study centers on two critical finetuning tasks in text-to-image synthesis: subject-driven generation, which seeks to render images of a specific subject in novel contexts using limited subject imagery and text prompts; and controllable generation, in which the model integrates supplementary control inputs. Through empirical evaluation, we demonstrate that our OFT approach surpasses existing alternatives regarding image quality and training efficiency.
\end{abstract}

\section{Introduction}

Modern text-to-image diffusion models~\cite{saharia2022photorealistic,ramesh2022hierarchical,rombach2022high} deliver state-of-the-art performance for text-based, high-quality image generation. Nevertheless, reliance solely on textual prompts can introduce vagueness and may fall short in enabling nuanced or precise modulation over generated content. To confront this limitation, we focus on two principal variants of text-to-image generation tasks in this work:

\begin{itemize}[leftmargin=*,nosep]

    \item \textbf{Subject-driven generation}~\cite{ruiz2023dreambooth}: When provided with a small collection of images depicting a particular subject, this task involves synthesizing novel renditions of that subject situated in various contexts, guided by a textual prompt.
    \item \textbf{Controllable generation}~\cite{zhang2023adding,mou2023t2i}: In this setting, alongside a text prompt, the model receives an additional modality for control—such as canny edge maps or segmentation layouts—and is tasked with generating images conforming to these auxiliary instructions.
\end{itemize}

Both tasks fundamentally address the challenge of refining text-to-image diffusion models through finetuning, while maintaining the generative capabilities achieved during pretraining. The essential criteria for an ideal finetuning approach are: (1) \emph{training efficiency}, characterized by a reduced number of adjustable parameters and minimal training epochs, and (2) \emph{generalizability preservation}, referring to the retention of the model’s ability to produce diverse and high-quality outputs. In practice, two major strategies have emerged for finetuning: modifying neuron weights incrementally with a low learning rate (\emph{e.g.}, \cite{ruiz2023dreambooth}), or augmenting the network by incorporating a compact module with re-parameterized weights (\emph{e.g.}, \cite{hulora2022,zhang2023adding}). Although these methods are straightforward to implement, neither reliably ensures the conservation of the generative prowess gained during pretraining. Moreover, there remains an absence of systematic guidelines for crafting effective finetuning methodologies or optimizing hyperparameters like the duration of training. This challenge is largely due to the lack of robust metrics that gauge how much of the pretrained generative capacity is retained post-finetuning. Common finetuning protocols operate under the presumption that minimizing the Euclidean distance between the finetuned and pretrained models correlates with better retention of the original model’s abilities. Consequently, finetuning often employs an extremely small learning rate. However, while this rationale may hold under specific circumstances, we contend that Euclidean distance alone fails to encapsulate the extent of semantic fidelity preservation. Adopting a more structural metric to assess changes between the finetuned and pretrained models could significantly enhance both the retention of pretrained performance and the reliability of the finetuning process.

Motivated by empirical findings that hyperspherical similarity effectively captures semantic relationships~\cite{liu2017deep,liu2018decoupled,chen2020angular}, we adopt hyperspherical energy~\cite{liu2018learning} as a means to represent the pairwise relational arrangement among neurons. This concept quantifies the aggregate hyperspherical similarity (\emph{e.g.}, cosine similarity) across all neuron pairs within the same layer, reflecting the degree to which neurons are distributed uniformly on the unit hypersphere~\cite{liu2021learning}. We propose that an effective finetuned model should exhibit only minimal alterations in hyperspherical energy relative to the original pretrained model. A straightforward approach would involve incorporating a regularization term to stabilize hyperspherical energy throughout finetuning. However, this strategy does not ensure that the change in hyperspherical energy will be sufficiently minimized. To address this, we leverage a specific invariance characteristic of hyperspherical energy: the pairwise hyperspherical similarity between neurons remains exactly unchanged if all neurons undergo an identical orthogonal transformation. Building upon this principle, we introduce Orthogonal Finetuning~(OFT), a technique designed to adapt large text-to-image diffusion models for downstream tasks while maintaining constant hyperspherical energy. The main premise is to learn an orthogonal transformation, shared across a layer, that keeps the angles between neurons intact. In essence, OFT can be interpreted as modifying the reference coordinate frame for neurons in a given layer. Additionally, considering that a smaller Euclidean distance between the finetuned and the pretrained model generally preserves more pretraining characteristics, we develop an extension called Constrained Orthogonal Finetuning~(COFT), which restricts the updated model to lie within a fixed-radius hypersphere centered at the pretrained neuron positions.

The underlying rationale for employing orthogonal transformations in neuron finetuning partially derives from the principles of the 2D Fourier transform. Through this transform, an image is expressed in terms of its magnitude and phase spectra. Notably, the phase spectrum encodes angular relationships between the input and the basis functions, which are instrumental in retaining most of the semantic content. Empirical observations reveal that substituting an image’s magnitude spectrum with a random one while keeping its phase spectrum intact still allows for the reconstruction of the original image’s semantics. This observation points toward the centrality of modifying neuron orientations for semantic alteration of generated images—a central objective in both subject-driven and controllable image generation tasks. Nevertheless, granting unrestricted freedom in changing neuron directions can compromise the generative performance established during pretraining. To mitigate this risk, our approach seeks to maintain the inter-neuron angles, drawing on the premise that such angles are essential to the neural network's capacity to encode knowledge. Accordingly, we propose learning a shared orthogonal transformation across the neurons of each layer to preserve the original hyperspherical energy.

Our methodology is further informed by orthogonal over-parameterized training~\cite{liu2021orthogonal}, wherein neural networks intended for classification are trained from random initializations using orthogonal transformations. This approach stems from the fact that random initializations typically yield a low expected hyperspherical energy, which \cite{liu2021orthogonal} aims to maintain throughout training; such low energy correlates with superior classification generalization~\cite{liu2018learning,lin2020regularizing}. The findings of \cite{liu2021orthogonal} indicate that orthogonal transformations provide ample expressivity for constructing well-generalizing networks for classification tasks. By contrast, our focus lies in adapting text-to-image diffusion models for enhanced controllability and improved generative performance on downstream tasks. We note key distinctions between OFT and the approach presented in \cite{liu2021orthogonal}: (1) while the method of \cite{liu2021orthogonal} targets minimizing hyperspherical energy, our OFT framework is designed to conserve the hyperspherical energy from the pretrained model, thus maintaining the essential structure learned during pretraining—an imperative particularly for diffusion model finetuning, where energy minimization may disrupt critical semantic organization; (2) OFT explicitly seeks to limit divergence from the pretrained model, resulting in a constrained formulation, unlike the unconstrained setting of \cite{liu2021orthogonal}. The essence of effective finetuning lies in balancing adaptability with the preservation of pretrained capabilities, and we posit that our OFT approach successfully attains this equilibrium. Our key contributions are as follows:

\begin{itemize}[leftmargin=*,nosep]

    \item We introduce a new finetuning technique—Orthogonal Finetuning—designed to enhance controllability in text-to-image diffusion models. To address stability concerns, we also propose a constrained extension that bounds the angular divergence relative to the original pretrained parameters.
    \item In contrast to traditional finetuning methods, OFT enables model adaptation while guaranteeing the preservation of hyperspherical energy. Our empirical studies indicate that this energy is a critical factor for retaining the pretrained model's generative semantics.
    \item Our experiments involve applying OFT to subject-driven and controllable generation tasks. A comprehensive evaluation demonstrates that our approach consistently surpasses existing baselines, leading to better image quality, faster convergence, and more stable training dynamics. Additionally, OFT is notably sample-efficient, reaching satisfactory convergence with far fewer training samples and epochs.
    \item For controllable generation, we propose a novel control consistency metric to quantitatively assess model controllability. This metric is based on extracting the control signal from the generated image and comparing it to the original input signal, thereby offering additional support for the superior controllability of OFT.
\end{itemize}

\section{Related Work}

\textbf{Text-to-image diffusion models}. Recent years have witnessed remarkable advances~\cite{saharia2022photorealistic,ramesh2022hierarchical,rombach2022high,gu2022vector,nichol2022glide} in the field of text-to-image synthesis, spurred by innovations in diffusion-based generative modeling~\cite{ho2020denoising,song2021score,song2021denoising,dhariwal2021diffusion} and cross-modal representation learning~\cite{su2020vlbert,lu2019vilbert,radford2021learning,wang2021simvlm,li2022blip,li2023blip,alayrac2022flamingo,schuhmann2022laion}. Systems such as GLIDE~\cite{nichol2022glide} and Imagen~\cite{saharia2022photorealistic} train diffusion models directly in the pixel space. GLIDE employs joint training of the text encoder and diffusion model with paired datasets, while Imagen opts for a fixed, pretrained text encoder throughout. Alternatively, Stable Diffusion~\cite{rombach2022high} and DALL-E2~\cite{ramesh2022hierarchical} perform diffusion modeling within a compressed latent space: Stable Diffusion utilizes VQ-GAN~\cite{esser2021taming} for learning a visual codebook that serves as the latent representation, whereas DALL-E2 builds a joint latent space leveraging CLIP~\cite{radford2021learning} for unified image-text representations. Beyond diffusion models, approaches based on generative adversarial networks~\cite{reed2016generative,zhang2017stackgan,xu2018attngan,li2019controllable} and autoregressive architectures~\cite{ramesh2021zero,ding2021cogview,wu2022nuwa,yuscaling} have also been pursued for this task. OFT is designed to be agnostic to the underlying diffusion model architecture and can be broadly utilized across different text-to-image diffusion frameworks.

\textbf{Subject-driven generation}. Addressing the issue of unwanted alterations to the main subject, some methods~\cite{avrahami2022blended,nichol2022glide} introduce user-specified masks as auxiliary conditioning. Inversion-based techniques~\cite{choi2021ilvr,dhariwal2021diffusion,ramesh2022hierarchical,gal2022image} facilitate context editing while maintaining the subject's core features. Alternative approaches like~\cite{hertz2022prompt} achieve both local and global edits without requiring masks. However, these techniques often fall short in preserving fine, identity-specific characteristics. Pivotal Tuning~\cite{roich2022pivotal} mitigates this by adjusting the generator near an inverted latent code with added regularization for identity retention. Similarly, \cite{nitzan2022mystyle} focuses on capturing personalized generative priors from collections of facial images, while \cite{casanova2021instance} enables instance variation at the expense of losing unique details. DreamBooth~\cite{ruiz2023dreambooth} leverages a custom token and a handful of subject photos, optimizing the diffusion model using a reconstruction objective and a prior preservation loss for the subject's class. In our work, OFT leverages the DreamBooth pipeline, replacing standard finetuning procedures with orthogonal transformations instead of conventional low-rate parameter updates.

\textbf{Controllable generation}. Image-to-image translation represents a subtype of controllable generation, with prior research predominantly leveraging conditional generative adversarial networks~\cite{isola2017image,zhu2017toward,wang2018high,park2019semantic,choi2018stargan}. More recently, diffusion models have also been utilized for image-to-image translation tasks~\cite{saharia2022palette,voynov2022sketch,wang2022pretraining}. In particular, ControlNet~\cite{zhang2023adding} introduced a strategy for guiding pretrained diffusion models via finetuning in response to external control inputs, demonstrating notable success in controllability. Similarly, T2I-Adapter~\cite{mou2023t2i} enhances control over image generation through the finetuning of diffusion models. Consistent with the frameworks presented in \cite{zhang2023adding,mou2023t2i}, we implement OFT for finetuning diffusion models and observe that it delivers superior controllability, requiring fewer samples and fewer tunable parameters. Importantly, OFT does not incur any extra computation during inference.

\textbf{Model finetuning}. Tailoring large pretrained models for specific downstream tasks through finetuning has gained considerable traction in recent years~\cite{devlin2018bert,brown2020language,he2022masked}. Adaptation strategies, such as those explored in \cite{hulora2022,houlsby2019parameter,pfeiffer2020adapterfusion}, constitute a popular subset of finetuning, especially within the field of natural language processing. OFT is closely related to LoRA~\cite{hulora2022}, which employs low-rank structures for additive weight updates during finetuning. However, OFT updates neuron weights multiplicatively by applying a shared orthogonal transformation per layer. This approach guarantees preservation of the angles between neurons within a layer, resulting in improved stability.

\section{Orthogonal Finetuning}

\subsection{Why Does Orthogonal Transformation Make Sense?}

To motivate the use of orthogonal transformations in finetuning text-to-image diffusion models, we break down the discussion into two primary considerations: (1) the rationale for focusing on finetuning neuron angles (that is, their directions), and (2) the justification for utilizing orthogonal transformations in the adjustment of these angles.

To address the first question, we take cues from empirical studies such as \cite{liu2018decoupled,chen2020angular}, which identify that the difference in angular features effectively reflects the underlying semantic disparity. According to SphereNet~\cite{liu2017deep}, constraining all neural activations to lie on a unit hypersphere during training results in similar network capacity but often improved generalization, suggesting that neuron direction encodes the most salient data characteristics. To illustrate the significance of neuron angles, we present a simplified experiment shown in Figure~\ref{fig:toy_ae}, wherein a basic convolutional autoencoder is trained on flower image data. During training, feature maps are generated via the conventional inner product (with $z$ as the convolution kernel $\bm{w}$ applied to input $\bm{x}$ in the sliding window). For testing, we explore three distinct approaches to feature map computation: (a) inner product (as used in training), (b) magnitude-based, and (c) angle-based (angular information). The outcomes in Figure~\ref{fig:toy_ae} indicate that utilizing neuron angles alone suffices to almost completely reconstruct the original images, while relying on magnitude yields no discernible information. Notably, the cosine activation in (c) is absent during training, which relies exclusively on the inner product. These observations underscore the crucial role of neuron directions in encapsulating the semantic attributes of the input data. Consequently, adjusting neuron angles appears to be a promising strategy for manipulating the semantic content of images.

Addressing the second question, the most straightforward technique for refining neuron direction is to rotate or reflect all neurons within the same layer simultaneously, a process that essentially employs orthogonal transformation. While it is conceivable to adopt more granular angular manipulations by assigning individual rotation angles to each neuron, our experiments indicate that orthogonal transformations offer an optimal balance of expressiveness and regularization. Additionally, \cite{liu2021orthogonal} provides evidence that orthogonal transformation possesses sufficient representational power for neural network learning. To empirically validate this claim, we conduct an experiment—following the ControlNet~\cite{zhang2023adding} protocol—focused on controllable generation. As shown in Figure~\ref{fig:ortho}, our conventional OFT exhibits robust performance and precise controllability after completing training (by epoch 20). In stark contrast, a version of OFT omitting the orthogonality constraint is unable to produce realistic images or exert effective control, thereby confirming the critical necessity of the orthogonality requirement within OFT.

Traditionally, finetuning approaches rely on applying gradient descent with a reduced learning rate to adapt model parameters, either globally or in specific layers. By employing a small learning rate, such procedures inherently limit how much the updated model diverges from its original pretrained configuration. Consequently, standard finetuning can be viewed as approximately optimizing $\thickmuskip=2mu \medmuskip=2mu \| \bm{M} - \bm{M}^0\|$, where $\bm{M}$ represents the updated weights and $\bm{M}^0$ the pretrained ones. Although this constraint has clear heuristic justification, it might allow excessive freedom when adjusting large models. To further restrict the parameter update, LoRA augments this scheme with a low-rank requirement on the difference: $\thickmuskip=2mu \medmuskip=2mu \text{rank}(\bm{M} - \bm{M}^0)=r'$, for some small rank $r'$. In contrast, OFT imposes a constraint on the neuron-level relationships, specifically requiring $\thickmuskip=2mu \medmuskip=2mu \| \text{HE}(\bm{M})-\text{HE}(\bm{M}^0) \|=0$, where $\text{HE}(\cdot)$ computes the hyperspherical energy for a given weight matrix.

Consider, for illustration, a fully connected layer parameterized as $\thickmuskip=2mu \medmuskip=2mu \bm{W} = \{\bm{w}_1, ..., \bm{w}_n\} \in \mathbb{R}^{d \times n}$, with each column $\bm{w}_i \in \mathbb{R}^d$ corresponding to the $i$-th neuron (and $\bm{W}^0$ representing pretrained values). Given an input $\thickmuskip=2mu \medmuskip=2mu \bm{x} \in \mathbb{R}^d$, the output is $\thickmuskip=2mu \medmuskip=2mu \bm{z} = \bm{W}^\top \bm{x}$. Under OFT, optimization is framed as aligning the hyperspherical energies of $\bm{W}$ and $\bm{W}^0$:

\begin{equation}\label{eq:oft_principle}
\footnotesize
\min_{\bm{W}} \left\| \text{HE}(\bm{W})-\text{HE}(\bm{W}^0)\right\|~~\Leftrightarrow~~\min_{\bm{W}} \bigg{\|} \sum_{i\neq j} \|\hat{\bm{w}}_i-\hat{\bm{w}}_j\|^{-1}-\sum_{i\neq j} \|\hat{\bm{w}}^0_i-\hat{\bm{w}}^0_j\|^{-1}\bigg{\|}
\end{equation}

with the normalized neuron vectors denoted by $\thickmuskip=2mu \medmuskip=2mu \hat{\bm{w}}_i := \bm{w}_i / \| \bm{w}_i \|$, and the hyperspherical energy calculated as $\thickmuskip=2mu \medmuskip=2mu \text{HE}(\bm{W}) := \sum_{i\neq j} \|\hat{\bm{w}}_i-\hat{\bm{w}}_j\|^{-1}$. It is straightforward to note that Eq.~\eqref{eq:oft_principle} admits zero as its minimal value, which occurs whenever $\bm{W}$ and $\bm{W}^0$ differ only by an orthogonal transformation: $\thickmuskip=2mu \medmuskip=2mu \bm{W} = \bm{R}\bm{W}^0$ for orthogonal $\thickmuskip=2mu \medmuskip=2mu \bm{R} \in \mathbb{R}^{d \times d}$ (rotation if $\det(\bm{R}) = 1$, reflection if $\det(\bm{R}) = -1$). The central premise of OFT is precisely to limit adaptation to such transformations, where only orthogonal matrices—shared across each layer—are learned to manipulate the neurons. Formally, OFT seeks to adjust the orthogonal parameter $\thickmuskip=2mu \medmuskip=2mu \bm{R

\begin{equation}\label{eq:oft_general}
    \footnotesize
    \bm{z}=\bm{W}^\top\bm{x}=(\bm{R}\cdot\bm{W}^0)^\top\bm{x},~~~\text{s.t.}~\bm{R}^\top\bm{R}=\bm{R}\bm{R}^\top=\bm{I}
\end{equation}

Here, $\bm{W}$ represents the weight matrix obtained after OFT fine-tuning, while $\bm{I}$ stands for the identity matrix. An overview of OFT is given in Figure~\ref{fig:block}. In analogy to LoRA’s use of zero initialization, it is necessary for OFT to start fine-tuning the model exactly from the initial pretrained weights $\bm{W}^0$. To ensure this, the orthogonal matrix $\bm{R}$ is initialized as the identity matrix, allowing the fine-tuned parameters to match the pretrained weights at the outset. Maintaining the orthogonality constraint on $\bm{R}$ can be achieved via differentiable orthogonalization approaches, such as those proposed by \cite{liu2021orthogonal,lezcano2019cheap}. We further elaborate on how to uphold this orthogonality property while minimizing computational burden.

\subsection{Efficient Orthogonal Parameterization}\label{sect:eff_ortho}

Although classical orthogonalization procedures, such as the Gram-Schmidt process, can be made differentiable, their computational overhead makes them impractical for large-scale problems~\cite{liu2021orthogonal}. For improved computational tractability, we employ Cayley parameterization to construct the orthogonal matrix. In this scheme, $\bm{R}$ is generated as $\thickmuskip=2mu \medmuskip=2mu \bm{R}=(\bm{I}+\bm{Q})(I-\bm{Q})^{-1}$, where $\bm{Q}$ is a skew-symmetric matrix satisfying $\thickmuskip=2mu \medmuskip=2mu \bm{Q}=-\bm{Q}^\top$. This method enhances efficiency but imposes a minor limitation: Cayley parameterization yields orthogonal matrices with determinant $1$, corresponding to the special orthogonal group. Nevertheless, empirical observations suggest that this constraint has negligible practical effect on model performance. Even with Cayley parameterization, using a full $\bm{R}$ may incur significant parameter costs for large $d$. To alleviate this, we suggest modeling $\bm{R}$ as a block-diagonal matrix comprising $r$ blocks, each of dimension $\frac{d}{r}$, resulting in the structure:

\begin{equation}
    \footnotesize
    \bm{R}=\text{diag}(\bm{R}_1,\bm{R}_2,\cdots,\bm{R}_r)=\begin{bmatrix}
    \bm{R}_1\in \text{O}(\frac{d}{r}) &  &\\
    &\ddots & \\
    & & \bm{R}_r\in \text{O}(\frac{d}{r})
    \end{bmatrix}\in \text{O}(d)
\end{equation}

Here, $\text{O}(d)$ represents the orthogonal group in $d$ dimensions, with $\thickmuskip=2mu \medmuskip=2mu \bm{R}\in\mathbb{R}^{d\times d}$ and, for all $i$, $\thickmuskip=2mu \medmuskip=2mu \bm{R}_i\in\mathbb{R}^{d/r\times d/r}$ serving as orthogonal matrices. In the special case where $\thickmuskip=2mu \medmuskip=2mu r=1$, the block-diagonal orthogonal structure collapses to the standard full-rank unconstrained orthogonal matrix. The parameter count for a $d\times d$ orthogonal matrix is given by $\thickmuskip=2mu \medmuskip=2mu d(d-1)/2$, equating to a complexity of $\mathcal{O}(d^2)$. In contrast, for an orthogonal matrix with $r$ block diagonals, the total number of parameters is $\thickmuskip=2mu \medmuskip=2mu d(d/r-1)/2$, leading to a reduced parameter complexity of $\thickmuskip=2mu \medmuskip=2mu \mathcal{O}(d^2/r)$. An additional reduction is possible by sharing block matrices across all blocks, such that $\thickmuskip=2mu \medmuskip=2mu\bm{R}_i=\bm{R}_j$ for all $i\neq j$, yielding a parameter complexity of only $\mathcal{O}(d^2/r^2)$. Crucially, even with these parameter-saving modifications, the matrix $\bm{R}$ retains its orthogonal structure, and thus the preservation of hyperspherical energy is not compromised.

We next compare the parameter efficiency of OFT to that of LoRA. In LoRA, using a low-rank parameter $r'$, the number of tunable parameters is $\thickmuskip=2mu \medmuskip=2mu r'(d+n)$. Assuming both $r$ for OFT and $r'$ for LoRA scale with the model’s width $d$—for instance, $\thickmuskip=2mu \medmuskip=2mu r=r'=\alpha d$, with some constant $\thickmuskip=2mu \medmuskip=2mu 0<\alpha\leq 1$—the LoRA approach reaches a parameter complexity of $\mathcal{O}(d^2+dn)$, while OFT attains a more efficient $\mathcal{O}(d)$. To concretely illustrate the disparity in parameter complexity between OFT and LoRA, consider a scenario in which the weight matrix is of size $\thickmuskip=2mu \medmuskip=2mu 128\times 128$; for $\thickmuskip=2mu \medmuskip=2mu r'=8$, LoRA involves $2,048$ parameters, whereas OFT—without block sharing and with $\thickmuskip=2mu \medmuskip=2mu r=8$—only requires $960$ trainable parameters.

\subsection{Constrained Orthogonal Finetuning}

A further restriction of the standard OFT is possible by limiting how much the finetuned model may deviate from its original pretrained configuration. This leads to the Constrained Orthogonal Finetuning (COFT) variant, where the forward operation is as follows:

\begin{equation}\label{eq:coft_general}
    \footnotesize
    \bm{z}=\bm{W}^\top\bm{x}=(\bm{R}\cdot\bm{W}^0)^\top\bm{x},~~~\text{s.t.}~\bm{R}^\top\bm{R}=\bm{R}\bm{R}^\top=\bm{I},~\norm{\bm{R}-\bm{I}}\leq \epsilon
\end{equation}

This formulation introduces both the standard orthogonality constraint and an additional bound on the distance between $\bm{R}$ and the identity matrix, with a deviation tolerance of $\epsilon$. The Cayley parameterization, as detailed in Section~\ref{sect:eff_ortho}, can be used to satisfy the orthogonality requirement. Nonetheless, enforcing the $\epsilon$-bounded deviation via the Cayley parameterization proves challenging. To elucidate how the Cayley transform behaves, the Neumann series can be applied for the approximation $\

Since the series converges with respect to the operator norm, we are permitted to impose the constraint $\thickmuskip=2mu \medmuskip=2mu\|\bm{R}-\bm{I}\|\leq\epsilon$ within the Cayley transform framework. This can be equivalently rewritten as a constraint on the matrix $\bm{Q}$: $\thickmuskip=2mu \medmuskip=2mu \|\bm{Q}-\bm{0}\|\leq\epsilon'$, where $\bm{0}$ is the zero matrix and $\epsilon'$ serves as a separate, possibly distinct tolerance parameter from $\epsilon$. Enforcing this new restriction on $\bm{Q}$ can be efficiently handled via projected gradient descent. To ensure that the orthogonal matrix $\bm{R}$ begins as the identity, we initialize $\bm{Q}$ to the zero matrix. The COFT approach is characterized by imposing two explicit conditions: minimization of hyperspherical energy variation, and a bounded departure from the weights of the pretrained network. While the second requirement is generally only implicitly maintained in standard finetuning methods, COFT formalizes it as an explicit condition. Despite the strong results obtained with OFT, we observe that COFT offers enhanced stability during finetuning, thanks to its explicit constraint on model deviation. Figure~\ref{fig:eps_coft} demonstrates the influence of $\epsilon$ on COFT's performance: as $\epsilon$ increases, COFT increasingly approximates OFT behavior, whereas reducing $\epsilon$ leads the COFT-adapted model to more closely follow the original pretrained text-to-image diffusion model.

\subsection{Re-scaled Orthogonal Finetuning}

We introduce a straightforward extension to OFT by also learning a per-neuron magnitude scaling factor. The rationale behind this modification is that scaling the neuron outputs does not alter the hyperspherical energy, since normalization removes such effects. Concretely, our modified forward computation is: $\thickmuskip=2mu \medmuskip=2mu\bm{z}=(\bm{R}\bm{W}^0\bm{D})^\top\bm{x}$\footnote{\textbf{Errata}: The camera-ready NeurIPS submission incorrectly stated the re-scaled OFT forward equation as $\thickmuskip=2mu \medmuskip=2mu\bm{z}=(\bm{D}\bm{R}\bm{W}^0)^\top\bm{x}$, though the actual code is correct and results in Appendix~\ref{app:rescaled} remain unaffected.}, where $\thickmuskip=2mu \medmuskip=2mu \bm{D}=\text{diag}(s_1,\cdots,s_n)\in\mathbb{R}^{n\times n}$ is a diagonal matrix with trainable entries $s_1,\dots, s_n$, each strictly positive. Differing from the original OFT formulation in Eq.~\eqref{eq:oft_general}, which only treats $\bm{R}$ as learnable, this re-scaled version allows both $\bm{D}$ and $\bm{R}$ to be optimized. The re-scaled OFT increases modeling capacity while incurring a negligible increase in parameter count. In our experiments, we employ the original OFT variant to highlight the impact of orthogonal transformations alone, though we observe that the re-scaled version typically yields improved outcomes (details are provided in Appendix~\ref{app:rescaled}).

\section{Intriguing Insights and Discussions}

\textbf{OFT's versatility across neural architectures}. In theory, OFT is applicable to a broad range of neural network models. For example, when used with Transformers, LoRA is often deployed on attention layers~\cite{hulora2022}. To ensure fair comparison with LoRA, our experiments constrain OFT to only adjust attention layer parameters. Furthermore, OFT's compatibility extends to convolutional layers; due to the block-diagonal nature of $\bm{R}$, it admits interesting interpretations not available with LoRA. For instance, assigning a block in $\bm{R}$ to each input channel restricts transformations to individual neuron channels, resembling depth-wise convolution operations~\cite{chollet2017xception}. Alternatively, by tying all blocks in $\bm{R}$, the transformation applies the same orthogonal operation across all channels. A preliminary exploration of convolutional layer finetuning using OFT is reported in Appendix~\ref{app:convolution}.

\textbf{Connection to LoRA}. LoRA incorporates a low-rank matrix to ensure that the underlying pretrained weights do not deviate significantly. By contrast, OFT enforces orthogonality (full-rank constraint) on the transformation applied to pretrained weights, thereby preserving the information acquired during pretraining. The forward computation for OFT can be represented as $\thickmuskip=2mu \medmuskip=2mu \bm{z}=(\bm{R}\bm{W}^0)^\top\bm{x}=(\bm{W}^0+(\bm{R}-\bm{I})\bm{W}^0)^\top\bm{x}$, where the term $\thickmuskip=2mu \medmuskip=2mu (\bm{R}-\bm{I})\bm{W}^0$ serves a similar role to the low-rank modification in LoRA. Assuming $\bm{W}^0$ maintains full rank, OFT effectively executes a low-rank update when $\thickmuskip=2mu \medmuskip=2mu \bm{R}-\bm{I}$ itself is of low rank. Analogous to the rank hyperparameter $r'$ in LoRA, OFT introduces a block-diagonal size parameter $r$ to control parameter efficiency. Intriguingly, LoRA and OFT offer contrasting approaches for efficient parameterization: LoRA leverages low-rank approximations, whereas OFT focuses on parameter reduction through structural sparsity—specifically, block-diagonal orthogonal transformations.

\textbf{Why OFT converges faster}. As Figure~\ref{fig:toy_ae} suggests, altering the directions of neurons leads to the most substantial semantic changes, aligning with the mechanism employed by OFT. Additionally, OFT can be interpreted as adapting neuron parameters on the surface of a smooth hypersphere, resulting in more favorable optimization properties. This characteristic is empirically corroborated in \cite{liu2021orthogonal}.

\textbf{Why not minimize hyperspherical energy}. Unlike \cite{liu2021orthogonal}, our method does not focus on reducing hyperspherical energy. In the context of classification tasks, a lack of redundancy among neurons is preferable. Achieving minimal hyperspherical energy would spatially distribute neurons uniformly over the hypersphere, which does not necessarily benefit finetuning since it risks erasing the informative structure captured during pretraining.

\textbf{Balancing flexibility and regularization in finetuning}. Our investigation reveals an inherent compromise between flexibility and regularization during finetuning. Traditional finetuning achieves maximum adaptability but is prone to instability and can lead to model collapse. Remarkably, the straightforward OFT approach effectively strikes a balance, maintaining flexibility while introducing regularity through conservation of neuron pairwise angular relationships. The block-diagonal parameterization further serves to enforce stronger regularization on the orthogonal matrix.

\textbf{No extra cost at inference time}. Contrary to methods like ControlNet, the OFT scheme does not impose any additional computational load during inference. At test time, the optimized orthogonal matrix $\bm{R}$ can be directly incorporated into the original pretrained weight matrix $\bm{W}^0$ by matrix multiplication, yielding a new effective weight matrix $\thickmuskip=2mu \medmuskip=2mu \bm{W}=\bm{R}\bm{W}^0$. As a result, the inference efficiency matches that of the pretrained model.

\section{Experiments and Results}

\textbf{Experimental protocol}. All experiments utilize Stable Diffusion v1.5~\cite{rombach2022high} as the foundational text-to-image generator. For an unbiased evaluation, generated outputs are randomly sampled across different techniques. Subject-driven generation procedures primarily follow the guidelines of DreamBooth~\cite{ruiz2023dreambooth}, while controllable generation methods adhere to practices from ControlNet~\cite{zhang2023adding} and T2I-Adapter~\cite{mou2023t2i}. To ensure comparability with LoRA, OFT or COFT modifications are restricted to the identical layers targeted by LoRA. Supplemental findings and experimental details can be found in Appendix~\ref{app:settings}.

\subsection{Subject-driven Generation}

\textbf{Protocol}. DreamBooth~\cite{ruiz2023dreambooth} and LoRA~\cite{hulora2022} serve as reference baselines. Each method implements the same loss function as specified in DreamBooth. For both DreamBooth and LoRA, we replicate the methodology and employ the optimal hyperparameter settings as proposed in the original work. Additional experimental results are documented in Appendix~\ref{app:settings},\ref{app:coft_oft},\ref{app:more_qualitative},\ref{app:failure}.

\textbf{Finetuning stability and convergence}. We begin by assessing the stability and rate of convergence during finetuning for DreamBooth, LoRA, OFT, and COFT. Figure~\ref{fig:teaser} and Figure~\ref{fig:db_convergence} summarize the findings. The results indicate that both COFT and OFT facilitate stable finetuning of the diffusion model. By 400 iterations, DreamBooth and OFT variants demonstrate strong control, in contrast to LoRA, which struggles to maintain subject identity. After 2000 iterations, DreamBooth exhibits image collapse, and LoRA misattributes visual features (producing yellow fur instead of a yellow shirt). Notably, OFT and COFT maintain robust and reliable control throughout the process. These observations provide evidence for the rapid and stable convergence of the OFT framework in subject-driven generative tasks. Importantly, the methods’ robustness to the chosen number of finetuning iterations reduces the burden of tuning this hyperparameter across different subjects. For both OFT and COFT, practitioners can select a relatively large iteration count without extensive adjustment. With a suitable choice of $\epsilon$ for COFT, setting both the learning rate and the number of iterations becomes straightforward.

\textbf{Quantitative comparison}. Adopting the protocol from \cite{ruiz2023dreambooth}, we perform a quantitative evaluation of subject fidelity (DINO~\cite{caron2021emerging}, CLIP-I~\cite{radford2021learning}), text prompt fidelity (CLIP-T~\cite{radford2021learning}), and sample diversity (LPIPS~\cite{zhang2018unreasonable}). Specifically, CLIP-I measures the mean pairwise cosine similarity between CLIP embeddings for generated versus real images, while DINO applies the same metric using ViT S/16 DINO embeddings. CLIP-T calculates the mean cosine similarity between the CLIP embeddings of the input text and those of the generated images. Additionally, we compute the average LPIPS cosine similarity for images generated of the same subject from an identical text prompt. Table~\ref{table:dreambooth_final} demonstrates that both COFT and OFT significantly surpass DreamBooth and LoRA on DINO and CLIP-I metrics, and deliver marginally superior or similar results in terms of prompt fidelity and diversity. To minimize stochastic effects, each method is repeatedly finetuned using 30 random seeds on the same pretrained model. The compiled results highlight that OFT not only offers improved convergence and training stability but also leads to more consistent and superior overall outcomes.

\textbf{Qualitative comparison}. To facilitate a more intuitive grasp of the advantages offered by OFT, we present several randomly selected cases of subject-driven generation in Figure~\ref{fig:dreambooth_final}. To ensure fairness, all shown samples are created using the same finetuned model per method, and the text prompts remain unaltered and are not specifically optimized for individual outputs. For each approach, the version of the model that attained the highest validation CLIP scores was chosen. The visualization in Figure~\ref{fig:dreambooth_final} clearly indicates that OFT and COFT both excel at preserving the semantic identity of the subject, whereas LoRA frequently struggles to retain subject characteristics (\emph{e.g.}, LoRA completely fails to maintain the bowl's identity in one example). Additionally, OFT and COFT offer notably finer control through text guidance, whereas DreamBooth, despite its proficiency at identity preservation, often fails to render images consistent with the given prompts (\emph{e.g.}, the bowl instance in the first row). These qualitative findings confirm that the OFT approach yields a superior balance of controllability and subject retention. Notably, OFT's performance is robust across varying iteration counts, enabling successful generation even when the number of iterations is large—an aspect where DreamBooth and LoRA do not perform as reliably. Further qualitative results are included in Appendix~\ref{app:more_qualitative}. Additionally, our human evaluation detailed in Appendix~\ref{app:human} provides further evidence supporting OFT's superior performance.

\subsection{Controllable Generation}

\textbf{Settings}. For baseline comparisons, we adopt ControlNet~\cite{zhang2023adding}, T2I-Adapter~\cite{mou2023t2i}, and LoRA~\cite{hulora2022}. The primary paper investigates three complex controllable generation scenarios: Canny edge to image~(C2I) leveraging the COCO dataset~\cite{lin2014microsoft}, segmentation map to image~(S2I) utilizing the ADE20K dataset~\cite{zhou2017scene}, and landmark to face~(L2F) on the CelebA-HQ dataset~\cite{karras2018progressive,xia2021tedigan}. All listed methods involve finetuning Stable Diffusion~(SD) v1.5 for 20 epochs on each of these datasets. Expanded results and additional experiments are available in Appendix~\ref{app:more_qualitative},~\ref{app:more_control_task},~\ref{app:failure}.

\textbf{Convergence}. To assess the convergence rate of ControlNet, T2I-Adapter, LoRA, and COFT on the L2F task, we conduct both quantitative and qualitative analyses. Quantitatively, we utilize the mean $\ell_2$ distance between predicted and control face landmarks as the primary metric. Figure~\ref{fig:face_convergence} displays the evolution of face landmark error across different finetuning epochs for each model. The data reveal that COFT and OFT substantially outpace others in terms of convergence speed: notably, LoRA requires 20 epochs to match the accuracy that OFT achieves at just 8 epochs. It is important to highlight that both OFT and COFT maintain a comparable parameter count to LoRA (and considerably fewer than ControlNet) while achieving much faster convergence than prior approaches. Additional evidence for OFT’s rapid convergence is seen in Figure~\ref{fig:teaser}; here, the rightmost example illustrates that OFT exhibits superior data efficiency compared to ControlNet and LoRA by converging effectively with merely 5\% of the ADE20K dataset. For the qualitative comparison, we concentrate on OFT, COFT, and ControlNet, given that ControlNet approaches OFT in landmark error performance. As seen in Figure~\ref{fig:face_convergence_qual}, both OFT and COFT consistently converge with the synthesized face poses progressively conforming to the control landmarks. Unlike our approaches, ControlNet’s controllability materializes abruptly after the 8th epoch—a convergence pattern previously documented in \cite{zhang2023adding}. This greater stability in convergence exhibited by OFT translates to smoother, more predictable training behavior, thereby facilitating hyperparameter tuning and making OFT particularly practical for real-world applications.

\textbf{Quantitative comparison}. To assess controllable generation, we propose a control consistency metric that measures how closely the control signal derived from the output image matches the original input control signal. In the context of the C2I task, Intersection-over-Union (IoU) and F1 score are utilized as evaluation metrics. For the S2I task, we report both mean IoU and mean/overall accuracy, whereas, for the L2F task, we determine the average $\ell_2$ distance between the ground truth control landmarks and those predicted. Detailed descriptions of these consistency metrics can be found in Appendix~\ref{app:settings}. Every competing method is evaluated using its optimal hyperparameters. The quantitative results in Table~\ref{table:control_gen} indicate that OFT and COFT offer notably superior and more precise control compared to alternative approaches. We note that methods based on adapters (such as T2I-Adapter and ControlNet) exhibit slower convergence and ultimately inferior outcomes. Specifically, LoRA surpasses ControlNet on S2I but is outperformed by ControlNet in C2I and L2F settings. Broadly, even with carefully selected hyperparameters, the maximum achievable performance with LoRA remains limited. In contrast, our OFT framework demonstrates ongoing improvements as the trainable parameter count increases, with no evidence of reaching a plateau. Importantly, our proposed framework for quantitative evaluation of controllable generation stands as a key contribution, offering an effective assessment of fine-tuned models’ control accuracy (subject to the precision of pre-trained segmentation or detection models).

\textbf{Qualitative comparison}. In addition to quantitative metrics, we present a qualitative assessment of OFT and COFT against leading approaches such as ControlNet, T2I-Adapter, and LoRA. The sample images displayed in Figure~\ref{fig:control_final} illustrate that OFT and COFT are capable of producing not only realistic and high-quality images, but also achieving precise control in generation. Specifically, for the S2I task, LoRA fails to adhere to the input segmentation map, whereas ControlNet, OFT, and COFT succeed in maintaining appropriate control over the generated outputs. OFT and COFT surpass ControlNet by synthesizing images with finer details, more lifelike textures, and improved geometric coherence, all while utilizing significantly fewer model parameters. In the C2I task, both OFT and COFT are adept at generating semantically faithful images from sparse Canny edge maps, a task at which T2I-Adapter and LoRA exhibit inferior performance. Regarding the L2F task, our methods achieve the most precise pose alignment in generated face images, even for faces in challenging orientations. Across the S2I, C2I, and L2F scenarios, both OFT and COFT demonstrate qualitative superiority over other state-of-the-art techniques, thereby highlighting the efficacy of the OFT framework for controllable image generation. For a broader set of qualitative results, additional examples for all three control tasks can be found in Appendix~\ref{app:control_exp}. Furthermore, Appendix~\ref{app:more_control_task} provides evidence that OFT generalizes effectively to a wider range of control scenarios, including dense pose-to-human body, sketch-to-image, and depth-to-image tasks.

\section{Concluding Remarks and Open Problems}

Driven by insights regarding the central role of angular relationships among neurons in shaping visual semantics, we introduce a straightforward yet powerful finetuning strategy—orthogonal finetuning (OFT)—for exerting control over text-to-image diffusion models. Our method is specifically applied to two important scenarios: subject-driven generation and controllable generation. In comparison with current approaches, OFT achieves superior control and stable finetuning while requiring fewer adjustable parameters. Crucially, OFT imposes no additional computational cost during inference, thus enabling efficient deployment.

OFT gives rise to several notable open questions. Firstly, orthogonality within OFT is ensured via Cayley parametrization, which necessitates computing a matrix inverse—this requirement places constraints on OFT’s scalability. Although the adoption of block diagonal parametrization offers a partial solution, developing faster, differentiable matrix inversion techniques is still an unresolved issue. Secondly, OFT possesses intriguing properties concerning compositionality: the orthogonal matrices derived from different OFT finetuning tasks can be composed through multiplication and retain orthogonality. Whether this property enables retention of knowledge across multiple downstream tasks is a promising avenue for further investigation. Finally, while OFT is parameter-efficient due to the use of block diagonal structure, this design choice inevitably leads to certain biases and curtails model flexibility. Identifying methods to enhance parameter efficiency with fewer biases represents a key unresolved challenge.